\documentclass[11pt]{article}
\usepackage[margin=1in]{geometry}
\usepackage[T1]{fontenc}
\usepackage{lmodern,microtype}
\usepackage{amsmath,amssymb,amsthm,booktabs,longtable,hyperref}
\newtheorem{lemma}{Lemma}
\newtheorem{theorem}{Theorem}
\newcommand{\pair}[2]{\pi(#1,#2)}
\newcommand{\tri}{\operatorname{tri}}
\title{The radix-independent formula constructor and its concrete source binding}
\author{Supplement to the PA194 radix-5 and radix-6 construction}
\date{13 November 2026, Pacific time}
\setlength{\emergencystretch}{2em}
\begin{document}
\maketitle

This note supplies the arithmetic macro argument needed to pass from the
literal primitive-instruction contracts to the new formula codes.  It is
a mathematical proof, with executable checks of all concrete row bindings.
It is not presented as a new Lean theorem.  The inherited source modules
\texttt{M6/Source/AXCons.lean}, \texttt{AXPair.lean}, and
\texttt{AXTimes4.lean} prove the corresponding radix-4 facts; the argument
below makes the dependence on the multiplier explicit.

\section{Registers, pairing, and exact source rows}

Write $T=r_1$, $W=r_2$, $X=r_6$, and $Y=r_7$.  All other registers form an
arbitrary frame $f$ that the macro preserves.  Set
\[
 \tri(s)=s(s+1)/2,\qquad \pair ab=\tri(a+b)+b,
 \qquad (\operatorname{hd}(w),\operatorname{tl}(w))=\pi^{-1}(w).
\]
The stack continues to use this unscaled Cantor pairing.  Only formula
implication changes to $I_k(a,b)=k\pair ab$.

Each occurrence of \texttt{node65} has the following resolved source
instructions, with all target addresses relative to its entry address $c$.
For a decrement, the two targets are the zero and positive targets,
respectively.  A positive decrement subtracts one.  An increment adds one.

\begin{center}
\begin{tabular}{rll@{\qquad}rll}
\toprule
Offset & Operation & Target(s) & Offset & Operation & Target(s)\\
\midrule
0 & decrement $W$ & $8,1$ & 11 & decrement $T$ & $15,12$\\
1 & increment $Y$ & $2$ & 12 & increment $Y$ & $13$\\
2 & decrement $X$ & $4,3$ & 13 & increment $X$ & $14$\\
3 & jump & $7$ & 14 & jump & $11$\\
4 & decrement $Y$ & $7,5$ & 15 & decrement $X$ & $21,16$\\
5 & increment $X$ & $6$ & 16 & increment $Y$ & $17$\\
6 & jump & $4$ & 17 & decrement $X$ & $20,18$\\
7 & jump & $0$ & 18 & increment $T$ & $19$\\
8 & decrement $Y$ & $11,9$ & 19 & jump & $17$\\
9 & increment $W$ & $10$ & 20 & jump & $11$\\
10 & jump & $8$ & 21 & decrement $Y$ & $k+23,22$\\
\bottomrule
\end{tabular}
\end{center}
Offsets $22,\ldots,k+21$ increment $T$ and continue to the next offset;
offset $k+22$ jumps to $21$.  The exit is $c+k+23$.

The builder verifies this entire 21-instruction prefix and the complete
multiplier loop separately at every occurrence.  It checks that all 61
occurrences of \texttt{node61} are exactly the multiplier bodies of 61
distinct \texttt{node65} occurrences: the relative tree address of each
body is $1110$.  There are no additional uses of the multiplier body.  The
sole direct reference to \texttt{node61} in the named source definitions is
in \texttt{node62}.  Each outside edge is preserved under the recorded
source-address relocation, and no outside edge enters an increment-body
interior.  These are finite facts about the entire source, not sample runs.

\section{The three arithmetic phases}

\begin{lemma}[Unpair prefix]
Let $(a,z)=\pi^{-1}(w)$.  Starting at $c$ with
$(T,W,X,Y)=(t,w,0,0)$, the first phase ends at $c+11$ with
$(t,z,a,0)$ and the same frame $f$.
\end{lemma}
\begin{proof}
At offset 0 one round consumes one unit of $W$.  If $X>0$, the round maps
$(X,Y)$ to $(X-1,Y+1)$ in five instructions.  If $X=0$, offsets 4--6
transfer the newly incremented $Y$ into $X$, so the round maps
$(0,Y)$ to $(Y+1,0)$ in $3Y+8$ instructions.  These are exactly the two
cases of the successor operation in the Cantor enumeration.  Induction
on the original value $w$ therefore reaches offset 8 with
$(W,X,Y)=(0,a,z)$.  Offsets 8--10 transfer $Y$ into $W$ in $3z+1$
instructions.  No instruction in this phase changes $T$ or $f$.
\end{proof}

For an exact duration, write
\[
 \mathcal{L}(0,p)=1,\qquad
 \mathcal{L}(n+1,p)=R(p)+\mathcal{L}(n,\operatorname{advance}(p)),
\]
where $R(a,b)=5$ for $a>0$, $R(0,b)=3b+8$, and
$\operatorname{advance}(a,b)=(a-1,b+1)$ for $a>0$ and
$(b+1,0)$ for $a=0$.  The prefix duration is
$\mathcal{L}(w,(0,0))+3z+1$.  A convenient bound is
$\mathcal{L}(w,(0,0))\leq w(3w+8)+1$, because every intermediate
component is at most $w$.

\begin{lemma}[Pair prefix]
Starting at $c+11$ with $(T,W,X,Y)=(t,z,a,0)$, the second phase ends at
$c+21$ with $(0,z,0,\pair at)$ and the same frame $f$.
\end{lemma}
\begin{proof}
Offsets 11--14 transfer each unit of $T$ into one unit of both $X$ and $Y$.
After $4t+1$ instructions the configuration at offset 15 is
$(0,z,a+t,t)$.  At offset 15, with $X=s+1$, the macro first subtracts
one from $X$ and adds one to $Y$.  Offsets 17--19 transfer the remaining
$s$ units of $X$ into $T$; offset 20 returns to 11; offsets 11--14
transfer those $s$ units back into $X$ and add $s$ to $Y$.
Thus the round maps $(X,Y)=(s+1,y)$ to $(s,y+s+1)$ in $7s+5$
instructions.  Induction on $X$ gives final $Y=t+\tri(a+t)=\pair at$.
The last zero decrement at offset 15 exits to 21.  The exact duration is
\[
 P(t,a)=4t+2+\sum_{j=0}^{a+t-1}(7j+5).
\]
The sum is empty if $a+t=0$.  Neither $W$ nor $f$ is changed.
\end{proof}

\begin{lemma}[Generic multiplication loop]
For every natural $k$, starting at offset 21 with
$(T,W,X,Y)=(t,z,x,y)$, the loop reaches offset $k+23$ with
$(t+ky,z,x,0)$ in exactly $(k+2)y+1$ source instructions.  The frame $f$
is unchanged.
\end{lemma}
\begin{proof}
If $y=0$, one zero decrement exits.  If $y>0$, one positive decrement,
the $k$ increments, and the back-jump take $k+2$ instructions and map
$(T,Y)$ to $(T+k,Y-1)$.  Induction on $y$ proves both the result and the
duration.  Equivalently, at every return to offset 21 the quantity
$T+kY$ is invariant and the nonnegative counter $Y$ decreases strictly
until exit.  The proof uses $Y=r_7$, not $X=r_6$.
\end{proof}

\begin{theorem}[Generic \texttt{Cons}]
For every natural $k$, every $t,w$, and every unrelated frame $f$, the
macro at $c$ terminates at $c+k+23$ with
\[
 (T,W,X,Y)=(I_k(\operatorname{hd}(w),t),\operatorname{tl}(w),0,0)
\]
when entered with $(t,w,0,0)$ and frame $f$.  Its exact source duration is
\[
 \mathcal{L}(w,(0,0))+3z+1+P(t,a)+(k+2)\pair at+1,
 \qquad (a,z)=\pi^{-1}(w).
\]
\end{theorem}
\begin{proof}
Compose the three lemmas.  At the start of the multiplier loop $T=0$
and $Y=\pair at$.  The output is therefore $k\pair at$, as claimed.
All three phases preserve the same arbitrary frame.  The duration is
positive and primitive recursive, and its termination proof uses only
ordinary natural-number induction available in PA.
\end{proof}

The changes for $k=5,6$ consequently alter formula construction while
preserving the tuple enumerator, the binary stack representation, and the
scheduler's pairing and unpairing operations.  Modus ponens uses the same
\texttt{Cons} macro as the axiom emitter, so its comparison consistently
uses $I_k$.  In particular, a primitive-contract pass alone is not being
used to infer this macro fact: the argument above and the 61 concrete
bindings supply that link.

\section{Formula tags and the free-constructor check}

The source's literal additions of one, two, and three remain unchanged.
They define
\[
 M_k(a,b)=I_k(a,b)+1,\quad
 U_k(a,B)=I_k(B,a)+2,\quad
 E_k(a,b)=I_k(a,b)+3.
\]
Here $a$ is the variable label bound by $U_k(a,B)$.  The abstract AX
interpreter in \texttt{build.py} treats $I_k$ as an uninterpreted binary
constructor, successor as an uninterpreted unary constructor, and input
values as independent atoms.  It executes all 190 macro tokens of the
unchanged principal axiom emitter.  It checks syntactic equality of the
entire output stack and top against the prescribed six retained axiom
heads, the restricted propositional head, the diagonal
$\Gamma_k(a)=I_k(a,U_k(a,a))$, and the guarded pointed-separation head.
The stack contains precisely ten new entries before the final guard.

This check excludes a literal use of \texttt{node61} in the AX token word;
the only two-increment literal is the unchanged universal tag.  Hence
changing the multiplier body does not accidentally change a formula tag.
It is an exact algebraic check for arbitrary inputs and arbitrary $k$;
it does not evaluate truth or establish the separate logical calibration.

The source binding also checks that the old and new sources have the same
3,429 primitive sites outside the increment bodies, with all operations,
register arguments, and successor edges preserved under relocation.
Their total sizes are therefore
\[
 3673+61(k-4)=
 \begin{cases}3734,&k=5,\\3795,&k=6.\end{cases}
\]
For the auxiliary radix-4 reassociation, the complete resolved instruction
list is identical to the principal PA194 source.

\section{Why the physical count stays at 194}

Let $i$ denote the one-increment reader and $ii$ the already available
two-increment reader.  The old multiplier body has shape $(ii,ii)$.
The new shapes are
\[
 4R:(i,(i,ii)),\qquad 5:((i,ii),ii),\qquad 6:((i,ii),(i,ii)).
\]
All introduce one shared new reader, native state 330, for $(i,ii)$.
The source reassociation ensures that the old two-increment reader,
native state 99, is now entered only immediately after a transition that
writes 1 and moves right.  Native state 41 is entered only immediately
after a transition that writes 0 and moves right.  Both can therefore
share one retained state 99 that first reads the previous tape cell.
Its actions are
\[
 (99,0)\mapsto(0,L,30),\qquad(99,1)\mapsto(1,L,2035).
\]

The four two-cell simulations below give the exact old successor after
three new steps.  The intervening paths visit only the old head cell and
its left neighbor, restore the appropriate old written symbol, and end
one cell right of the old head.  Here $h$ is the bit in the left neighbor
and $b$ is the originally scanned bit.

\begin{center}
\begin{tabular}{rrrl}
\toprule
Old state & $h$ & $b$ & New state path\\
\midrule
41 & 0 & 0 & $99\to30\to45\to24$\\
41 & 0 & 1 & $99\to2035\to61\to25$\\
99 & 1 & 0 & $99\to30\to46\to61$\\
99 & 1 & 1 & $99\to2035\to2036\to61$\\
\bottomrule
\end{tabular}
\end{center}

The builder checks these four cases together with all 28 cases of the
seven inherited folds on the final simultaneous transition table.  It
also checks every incoming row, so the history-bit hypotheses hold for
all entries, rather than merely for the blank run.  There are 135 fixed
controls and 67 axiom-reader controls before folding.  The eight
disjoint folds leave $135+67-8=194$ working controls and 388 transitions.
The separate halt target is not counted as a working control.

The dense table has initial control 0, working controls $0,\ldots,193$,
and halt target 194.  The tape alphabet is $\{0,1\}$, the blank is 0,
and every transition moves left or right.  Blank-start initialization
takes 58,320 physical steps and produces the exact complete source frame,
which is checked as equality of the entire finite-support tape.

\end{document}
